\chapter{Selected Source Code}
\label{app:code}

The full codebase is in the repository given in
Appendix~\ref{app:reproducibility}. This appendix reproduces the four excerpts
on which the claims of Chapters~\ref{chap:method} and~\ref{chap:results} most
directly depend, so that a reader can check the argument against the
implementation without cloning anything. Each listing is quoted verbatim from
the file and line range named in its caption; nothing is paraphrased for
presentation.

\section{Quote Grounding}

Listing~\ref{lst:quote-grounding} is the check the central claim of this work
rests on. \texttt{quote\_is\_grounded} accepts a quote only when its word
sequence appears \emph{contiguously} in the cited evidence, after typographic
and punctuation normalisation. Contiguity is what rejects a quote assembled
from two fragments, which is the failure a syntactic search for ellipsis
characters passes (Section~\ref{sec:generation}). The byte-level rule is kept
alongside it and reported as a second figure throughout, rather than replaced.

\begin{lstlisting}[style=thesis,caption={The contiguity definition of a verbatim quote, with the byte-level rule retained beside it. Source: \texttt{src/aviation\_rag/generation/verification.py}, lines 60--90.},label={lst:quote-grounding}]
def _normalize(value: str) -> str:
    """Whitespace and case only -- the original, strictest rule."""
    return WHITESPACE.sub(" ", value).strip().casefold()


def _normalize_typography(value: str) -> str:
    value = unicodedata.normalize("NFKC", value)
    for source, target in _TYPOGRAPHY.items():
        value = value.replace(source, target)
    return value


def _normalize_for_matching(value: str) -> str:
    """Typography- and punctuation-insensitive; word order is preserved."""
    value = NON_WORD.sub(" ", _normalize_typography(value))
    return WHITESPACE.sub(" ", value).strip().casefold()


def quote_is_grounded(quote: str, evidence_text: str) -> bool:
    """Does the quote's word sequence appear contiguously in the evidence?"""
    return _normalize_for_matching(quote) in _normalize_for_matching(evidence_text)


def quote_is_verbatim(quote: str, evidence_text: str) -> bool:
    """Byte-level check under the original whitespace-and-case rule."""
    return _normalize(quote) in _normalize(evidence_text)


def contains_elision(quote: str) -> bool:
    """Does the quote mark omitted material with an ellipsis?"""
    return bool(ELLIPSIS.search(quote))
\end{lstlisting}

\section{Structural Verification}

Listing~\ref{lst:verify-loop} is the loop that decides whether a claim is
supported. Two properties are worth reading off the code directly. A support
whose \texttt{source\_id} is absent from \texttt{evidence\_by\_id} fails
immediately --- this is the check that caught the fabricated identifier
\texttt{EU\_376\_2014\_ART\_5\_\_C001} on \texttt{QA\_0017}, inside a response
the provider's strict structured-output mode had already accepted. And a claim
can lose its support without contributing to \texttt{quote\_errors}, which is
why Section~\ref{sec:generation} reads support coverage from the list of
supported claim identifiers rather than inferring it from the decision.

\begin{lstlisting}[style=thesis,caption={The support loop: supplied-source admissibility, contiguity, and the byte-level second opinion. Source: \texttt{src/aviation\_rag/generation/verification.py}, lines 105--132.},label={lst:verify-loop}]
            claim_id = claim.get("claim_id", "?")
            claim_supported = True
            for support in claim.get("supports", []):
                if not isinstance(support, dict):
                    claim_supported = False
                    continue
                source_id = support.get("source_id")
                quote = support.get("quote")
                if (
                    source_id not in evidence_by_id
                    or not isinstance(quote, str)
                    or not quote.strip()
                ):
                    claim_supported = False
                    continue
                evidence_text = evidence_by_id[source_id]["text"]
                if contains_elision(quote):
                    elided_quotes.append(f"{claim_id}: elided quote from {source_id}")
                if not quote_is_grounded(quote, evidence_text):
                    quote_errors.append(
                        f"{claim_id}: quote not found in {source_id}"
                    )
                    claim_supported = False
                if not quote_is_verbatim(quote, evidence_text):
                    strict_errors.append(
                        f"{claim_id}: quote not byte-verbatim in {source_id}"
                    )
            if claim_supported and claim.get("claim_id"):
\end{lstlisting}

\section{Citation Scoring}

The evaluator that replaces the historical scorer resolves references against
the corpus instead of matching strings. Listing~\ref{lst:cite-exact} is the
exact measure, and the reason 10 of the 98 frozen records fall outside its
denominator rather than scoring zero: a gold citation that no parser can
resolve is usually one no reader can resolve either, and that count is reported
as a finding in Section~\ref{sec:invisible}.

\begin{lstlisting}[style=thesis,caption={Exact citation matching, with unresolvable gold labels excluded rather than failed. Source: \texttt{src/aviation\_rag/evaluation/citations.py}, lines 314--328.},label={lst:cite-exact}]
def exact_citation_match(prediction: str, gold_citation: str) -> tuple[bool | None, list[str]]:
    """Is every atomic unit of the gold citation named in the prediction?

    Returns ``(None, [])`` when the gold citation contains no parseable unit --
    for instance ``GM to Reg. (EU) No 376/2014 [04]`` or a bare ``Paragraph 3``.
    Those records are excluded from the exact metric's denominator rather than
    silently scored zero, and their count is itself a finding: a citation label
    a parser cannot resolve is usually one a reader cannot resolve either.
    """
    gold_units = parse_citations(gold_citation)
    if not gold_units:
        return None, []
    predicted = {citation.signature() for citation in parse_citations(prediction)}
    missing = [str(unit) for unit in gold_units if unit.signature() not in predicted]
    return (not missing), missing
\end{lstlisting}

Listing~\ref{lst:cite-candidates} handles the ambiguity the benchmark itself
carries. Where a gold citation names a provision without naming which of the
four regulations it belongs to, every regulation holding that provision counts
as a candidate. A prediction is therefore never marked wrong for failing to
draw a distinction the gold label does not draw.

\begin{lstlisting}[style=thesis,caption={Ambiguous gold citations resolve to a candidate set, not to a default attribution. Source: \texttt{src/aviation\_rag/evaluation/citations.py}, lines 138--158.},label={lst:cite-candidates}]
    def resolve_candidates(self, citation: Citation) -> set[str]:
        """Every parent the citation could denote, given what it actually says.

        A citation that names its regulation denotes exactly one parent. A bare
        ``Article 2`` denotes any of the four regulations that have an Article 2,
        and this returns all of them. Scoring against the candidate set rather
        than against the default attribution means a prediction is never marked
        wrong for failing to make a distinction that the gold citation itself
        does not make.
        """
        if not citation.defaulted_regulation:
            parent = self.resolve(citation)
            return {parent} if parent else set()
        candidates = set()
        for regulation in REGULATION_PREFIX:
            parent = self.resolve(
                Citation(regulation, citation.kind, citation.number, citation.subs)
            )
            if parent:
                candidates.add(parent)
        return candidates
\end{lstlisting}

